\subsection{Event-Scheduling Scheme}
Questo approccio si concentra sugli eventi come entità centrali che modificano lo stato del sistema in istanti temporali discreti.

Il simulatore mantiene una lista centrale ordinata chiamata \textbf{Scheduled Event List} (\textbf{SEL}). Ogni elemento della lista è una coppia composta da un tipo di evento e il momento esatto in cui accadrà. Il motore di simulazione funziona estraendo l'evento \textit{più imminente}, portando il tempo globale della simulazione a quel valore ed eseguendo la routine associata a quello specifico evento per aggiornare lo stato del sistema (es. un arrivo o una partenza).

La logica è prettamente legata alle attivazioni temporali e risulta "frammentata". Non esiste una visione globale della vita di un oggetto; ci sono solo blocchi isolati di codice che si attivano quando scatta un determinato clock.

\subsection{Process-Oriented Scheduling Scheme}
Questo approccio sposta il focus dagli eventi alle \textbf{entità} (es. i clienti in una coda, i pezzi in una catena di montaggio o i pacchetti in una rete) che attraversano il sistema.

Un \textit{processo} descrive l'intera sequenza di eventi e intervalli temporali che una singola entità sperimenta durante il suo ciclo di vita all'interno del sistema. Ad esempio:
\begin{enumerate}
	\item Arriva l'entità;
	\item Si controlla che il server è libero;
	\item Si riserva la risorsa;
	\item Attesa per la durata del servizio;
	\item Rilascio della risorsa;
\end{enumerate}

Nella simulazione orientata ai processi, ogni entità si comporta come se fosse un oggetto o un thread indipendente con la propria linea di controllo. Il codice viene scritto dal punto di vista dell'entità, e contiene funzioni che possono "sospendere" temporaneamente l'esecuzione del thread (es. funzioni di time delay o attese condizionate in coda prima di ottenere una risorsa).

\subsubsection{Summary schemes}
Molti linguaggi e software di simulazione moderni (come Arena o SIMSCRIPT citati nel testo) offrono un'interfaccia process-oriented perché è molto più intuitiva per il programmatore. Tuttavia, "sotto il cofano" (a livello di kernel software), spesso traducono tutto in un motore event-scheduling con una SEL centrale, perché è computazionalmente molto più efficiente da eseguire.

\begin{table}[H]
	\centering
	\label{tab:scheduling_schemes}
	\begin{tabularx}{\textwidth}{l X X}
		\toprule
		\textbf{Caratteristica} & \textbf{Event-Scheduling Scheme} & \textbf{Process-Oriented Scheme} \\
		\midrule
		\textbf{Focus principale} & Gli eventi atomici e isolati (Cosa succede al tempo $t$?) & Le entità e il loro ciclo di vita (Cosa fa questo oggetto nel sistema?) \\
		\addlinespace
		\textbf{Punto di vista del codice} & Frammentato in routine basate sul tipo di evento & Sequenziale e lineare per l'entità (simile alla logica di un thread) \\
		\addlinespace
		\textbf{Gestione del tempo} & Salti espliciti basati sulla \emph{Scheduled Event List} (SEL) & Sospensione del processo tramite funzioni di ritardo (\emph{delay}) \\
		\addlinespace
		\textbf{Flessibilità e utilizzo} & Universale, permette di modellare dinamiche molto complesse e arbitrarie & Ideale per sistemi di accodamento e logistica; richiede molte meno righe di codice \\
		\bottomrule
	\end{tabularx}
	
	\caption{Confronto tra Event-Scheduling Scheme e Process-Oriented Scheme.}
\end{table}

\subsection{Colored Petri Net (CPN)}

\subsubsection{Multiset}
Il formalismo matematico alla base dei colore nelle CPN si fonda sui multiset. Un multiset (o multiset non negativo) $\alpha$ definito su un insieme $D$ è una mappatura
\begin{equation}
	\alpha : D \rightarrow \mathbb{Z}
\end{equation}
Oppure riscrivibile come $\alpha : D \rightarrow \mathbb{N}$.

Viene rappresentato formalmente come la sommatoria degli elementi con peso non nullo:
\begin{equation}
	\alpha = \sum_{d\in D} \alpha(d) \oplus d
\end{equation}
In forma vettoriale, corrisponde a un vettore $\alpha =[\alpha(d_1)\alpha(d_2)\cdots \alpha(d_k)]^T \in \mathbb{Z}^k $.\newline

\noindent Operazioni sui multiset:
\begin{itemize}
	\item Somma ($\gamma = \alpha + \beta$): $\forall d \in D: \gamma (d) = \alpha (d) + \beta(d)$;
	\item Differenza ($\gamma = \alpha - \beta$): $\forall d \in D: \gamma (d) = \alpha (d) - \beta(d)$, può risultare negativa anche partendo da multiset non negativi;
	\item Prodotto per scalare ($\gamma = a\alpha$ con $a\in \mathbb{Z}$): $\forall d \in D: \gamma (d) = a\alpha (d)$;
	\item Confronto ($\alpha \le \beta$): $\forall d \in D: \alpha (d) \le \beta(d)$.
\end{itemize}

\subsubsection{Definizione generale di CPN}
Un rete di petri colorata è definita formalmente come una sestupla:
\begin{equation}
	N = (P, T, \text{Pre}, \text{Post}, Cl, Co)
\end{equation}
Dove:
\begin{itemize}
	\item $P$ è l'insieme degli $m$ posti;
	\item $T$ è l'insieme delle $n$ transizioni;
	\item $Cl$ è l'insieme globale dei colori ammissibili;
	\item $Co: P \cup T \rightarrow Cl$ è la \textbf{funzione colore} che associa a ciascun elemento un insieme ordinato e non vuoto di colori ricavati da $Cl$;
	\begin{itemize}
		\item Per un posto $p_i$, $Co(p_i)$ indica i colori possibili per i gettoni (token) contenuti in esso e $u_i$ rappresenta il loro numero;
		\item Per una transizione $t_j$, $Co(t_j)$ indica i possibili colori di occorrenza (o scatto) della transizione e $v_j$ rappresenta il loro numero;
	\end{itemize}
	\item $\text{Pre}$ e $\text{Post}$ sono le matrici di incidenza composte da mappature tra i colori dei posti e i multiset dei colori delle transizioni:
	\begin{equation}
		\text{Pre}(p,t) : Co(p)\rightarrow \mathcal{M}(Co(t))
	\end{equation}
	\begin{equation}
		\text{Post}(p,t) : Co(t)\rightarrow \mathcal{M}(Co(p))
	\end{equation}
	La matrice d'incidenza globale è definita come:
	\begin{equation}
		C = \text{Post} - \text{Pre}
	\end{equation}
\end{itemize}

\subsubsection{Marcatura e stato nelle CPN}
La marcatura $m_i$ di un posto $p_i$ è un multiset non negativo definito su $Co(p_i)$. La funzione associa a ogni potenziale colore di gettone un intero non negativo che indica quanti gettoni di quel colore sono presenti nel posto.

In formato vettoriale colonna di dimensione $u_i$, la componente $h$-esima $m_i(h)$ indica il numero di gettoni di colore $a_{i,h}$ nel posto $p_i$.

La marcatura complessiva delle CPN è un vettore colonna a $m$ dimensioni di multiset:
\begin{equation}
	m = [m_1 \cdots m_m]^T
\end{equation}

\noindent \textbf{Formalismo A: Guardie ed Espressioni sui nodi}
Questo approccio rende i modelli più leggibili dal punto di vista grafico.

\begin{itemize}
	\item \textbf{Funzione colore e transizioni}: la funzione $Co: P \rightarrow Cl$ mappa solo i posti. I colori delle transizioni non sono definiti esplicitamente ma implicitamente tramite funzioni di controllo logico dette \textbf{guardie} ($G(t)$), a valore booleano;
	\item \textbf{Binding}: un assegnamento di valori alle variabili delle espressioni tale per cui $G(t) = 1$ (quindi quando è vera) è chiamato \textit{binding}. L'insieme di tutti i binding possibili per $t$ è denotato come $B(t)$;
	\item Le matrici $\text{Pre}$ e $\text{Post}$ contengono espressioni le cui variabili sono i colori del posto connesso. Quando si applica un binding $\beta$, i valori della marcatura rimossa o aggiunta sono indicati come
	\begin{equation}
		\text{Pre}(\cdot, t)[\beta],
		\text{Post}(\cdot, t)[\beta]
	\end{equation}
	\item \textbf{Abilitazione e scatto}: una transizione $t$ è abilitata per un binding $\beta$ se e solo se la marcatura attuale soddisfa $m \ge \text{Pre}(\cdot, t)[\beta]$. Lo scatto produce una nuova marcatura
	\begin{equation}
		m' = m + \text{Post}(\cdot, t)[\beta] - \text{Pre}(\cdot, t)[\beta] = m + C(\cdot, t)[\beta]
	\end{equation}
\end{itemize}

\noindent \textbf{Formalismo B: Rappresentazione matriciale}
Questo approccio predilige l'algebra lineare ed è particolarmente utile per scopi di ricerca nel controllo di supervisione.
\begin{itemize}
	\item Le entrate delle matrici $\text{Pre}$ e $\text{Post}$ sono interamente costituite da matrici di numeri interi;
	\item $\text{Pre}(p_i, t_j)$ e $\text{Post}(p_i, t_j)$ diventano matrici di dimensioni $u_i \times v_j$ di interi non negativi. Ad esempio, l'elemento generico $\text{Pre}(p_i, t_j)(h, k)$ indica il peso dell'arco dal posto $p_i$ (rispetto al colore del gettone $a_{i,h})$) alla transizione $t_j$ (rispetto al colore di occorrenza $b_{j,k}$);
	\item \textbf{Abilitazione e scatto}: una transizione $t_j$ è abilitata rispetto al colore di occorrenza $b_{j,k}$ alla marcatura $m$ se e solo se $m_i(h) \ge \text{Pre}(p_i, t_j)(h,k)$ per ogni posto $p_i$ e per tutti gli $h= 1,\cdots, u_i$. Lo scatto aggiorna la marcatura per ogni posto secondo l'equazione lineare
	\begin{equation}
		m'_i(h) = m_i(h) + \text{Pre}(p_i, t_j)(h,k) - \text{Post}(p_i, t_j)(h,k) 
	\end{equation}
\end{itemize}

\subsubsection{Controllo di Supervisione e Gestione dei Conflitti}
\begin{itemize}
	\item \textbf{Transizioni controllabili e non controllabili}: l'insieme delle transizioni (o dei singoli colori di occorrenza) è diviso in due sottoinsiemi disgiunti.
	\begin{itemize}
		\item $T_c$, transizioni controllabili che il supervisore può abilitare o disabilitare tramite una variabile booleana $CV$;
		\item $T_{uc}$ transizioni non controllabili, governate esclusivamente dall'impianto;
	\end{itemize}
	Nel formalismo matriciale, un supervisore non può inibire una transizione non controllabile rispetto a un determinato colore;
	\item \textbf{Specifiche di stato (GMEC)}: i vincoli di sicurezza sono spesso espressi tramite Vincoli di Esclusione Mutua Generalizzata (GMEC) che impongono restrizioni sulle marcature legali ammissibili del sistema (es. limitare la capacità dei buffer dei robot o delle macchine nei sistemi flessibili di produzione);
	\item \textbf{Risoluzione dei conflitto}: due transizioni sono in conflitto effettivo se sono entrambe abilitate alla marcatura attuale, ma la quantità di gettoni disponibili non permette lo scatto simultaneo di entrambe. Per rendere deterministico l'avanzamento temporale e logico, la teoria prevede politiche di schedulazione basate su:
	\begin{itemize}
		\item Ordine di priorità statico/predefinito;
		\item Frequenze/probabilità di scatto identiche o stocastiche (routing rates);
		\item Algoritmi o logiche esterne guidate da supervisori basati sullo stato;
	\end{itemize}
	\item \textbf{Gestione del tempo (CTPN)}: le CPN temporizzate introducono una funzione tempo $f$ che associa a ogni binding o colore di occorrenza un ritardo deterministico o stocastico tra l'abilitazione e l'effettivo scatto atomico, permettendo valutazioni di performance (come il calcolo del throughput) attraverso semantiche a servente infinito (infinite server semantic).
\end{itemize}